
%\RequirePackage{pdf15}

\documentclass{beamer}

\usepackage[utf8]{inputenc}

\usepackage{mystyle}

\usepackage{tikz}
\usepackage{pgfplots}
\usepackage{subcaption}

%\usepackage{natbib}
\usepackage[style=authortitle,backend=biber]{biblatex}
\addbibresource{anthology.bib}
\addbibresource{emnlp2020.bib}
\renewcommand{\footnotesize}{\scriptsize}

\usepackage{tikz-dependency}
\usetikzlibrary{shapes.arrows, positioning, fit, bayesnet,
    arrows,backgrounds,patterns,matrix,calc,shadows,plotmarks,
    shapes,positioning,automata,positioning,spy,scopes,chains,decorations,decorations.pathreplacing}

\newcommand{\FancyUpArrow}{\begin{tikzpicture}[baseline=-0.3em]
\node[single arrow,draw,rotate=90,single arrow head extend=0.2em,inner
ysep=0.2em,transform shape,line width=0.05em,top color=green,bottom color=green!50!black] (X){};
\end{tikzpicture}}
\newcommand{\FancyDownArrow}{\begin{tikzpicture}[baseline=-0.3em]
\node[single arrow,draw,rotate=-90,single arrow head extend=0.2em,inner
ysep=0.2em,transform shape,line width=0.05em,top color=red,bottom color=red!50!black] (X){};
\end{tikzpicture}}

\AtBeginSection[]{
  \begin{frame}
  \vfill
  \centering
  \begin{beamercolorbox}[sep=8pt,center,shadow=true,rounded=true]{title}
    \usebeamerfont{title}\insertsectionhead\par%
  \end{beamercolorbox}
  \vfill
  \end{frame}
}

% quotes
\usepackage[style=british]{csquotes}

\def\signed #1{{\leavevmode\unskip\nobreak\hfil\penalty50\hskip1em
  \hbox{}\nobreak\hfill #1%
  \parfillskip=0pt \finalhyphendemerits=0 \endgraf}}

\newsavebox\mybox
\newenvironment{aquote}[1]
  {\savebox\mybox{#1}\begin{quote}\openautoquote\hspace*{-.7ex}}
  {\unskip\closeautoquote\vspace*{1mm}\signed{\usebox\mybox}\end{quote}}

%Information to be included in the title page:
\title{Word Games}
\author{J Chiu}

\setbeamertemplate{navigation symbols}{} 
\setbeamertemplate{footline}[frame number]

\begin{document}

\begin{frame}[plain]
\titlepage
\end{frame}

\begin{frame}
\frametitle{Language Games}
\begin{itemize}
\item Interaction is important 
    \begin{itemize}
    \item Cannot fully automate every task,
        i.e. task-oriented or information seeking dialogues
    \item Must handle diverse human input, although they may map
        to a low-dimensional manifold
    \end{itemize}
\item Interaction is hard 
    \begin{itemize}
    \item Human input is expensive and therefore limited
    \item In order to make certain problem aspects tractable,
        must make sacrifices in other
        (toy domain = out of distribution for pretrained models)
    \end{itemize}
\item What are the main challenges in language games,
    and what are the tradeoffs of different approaches?
\end{itemize}
\end{frame}

\begin{frame}
\frametitle{3 Challenges in Language Games}
\begin{enumerate}
\item Meaning representations
    \begin{itemize}
    \item A
    \end{itemize}
\item Search
    \begin{itemize}
    \item A
    \end{itemize}
\item Modeling other agents
    \begin{itemize}
    \item A
    \end{itemize}
\end{enumerate}
\end{frame}

\begin{frame}
\frametitle{Types of Games}
\begin{itemize}
\item 
\end{itemize}
\end{frame}

\begin{frame}
\frametitle{Axes of Tractability}
\begin{itemize}
\item Horizon length
    \begin{itemize}
    \item B
    \end{itemize}
\item Action space
    \begin{itemize}
    \item B
    \end{itemize}
\item A
    \begin{itemize}
    \item B
    \end{itemize}
\end{itemize}
\end{frame}

\begin{frame}
\frametitle{Meaning reps}
\begin{itemize}
\item In full generality, this problem is often encountered in hierarchical RL
    \begin{itemize}
    \item Less bleak in the language gamesetting
    \item Games are often very simple and can be constrained to small horizons,
        for example He He engineered a parser and policy that basically solves
        the negotiation task
    \end{itemize}
\item Many text-specific meaning representations (MR) to choose from
    \begin{itemize}
    \item Many are too complex
    \item Can we leverage existing MRs to learn a minimal task-specific
        representation that balances utility and expressivity?
    \end{itemize}
\end{itemize}
\end{frame}

\begin{frame}
\frametitle{Learning to plan}
\begin{itemize}
\item 
\end{itemize}
\end{frame}

\begin{frame}[allowframebreaks]
\frametitle{Citations}
\printbibliography
\end{frame}

\end{document}
